% TOP_PROBLEM: {"id":30007573,"problem_number":"LOCAL-30007573","title":"Voluntary participation over time","category_id":21,"category":{"id":21,"name":"economics","display_name":"Economics","description":"Open economic theory statements, published research agendas, and proposed scoped specifications; question provenance and historical priority are qualified per record.","slug":"economics","order_index":21,"created_at":"2026-10-08T00:00:00Z"},"status":"research_question","external_url":null,"legacy_ids":[],"published":false,"statement":"Characterize optimal dynamic contracts under costless sequential exit and finite liquidity, and quantify how participation constraints or limited escrow reduce revenue.","economics":{"original_id":62,"field":"Mechanism design and market design","kind":"theoretical","formalization_basis":"adapted_research_question","source_status":"research_agenda","priority_status":"earliest_found","priority_note":"This is the earliest source in which the relevant agenda was verified during this audit; absolute priority has not been established.","earliest_verified_reference":{"authors":["Dirk Bergemann","Juuso Välimäki"],"title":"Dynamic Mechanism Design: An Introduction","year":2018,"url":"https://cowles.yale.edu/sites/default/files/2022-09/d2102-r.pdf","doi":"10.1257/jel.20180892","locator":"§8 Concluding Remarks, pp.49–51 of June19,2018 revision","evidence_summary":"The conclusion explicitly identifies weaker common-prior assumptions, sequential participation, outside resale, and limited adoption of long-term contracts as research questions."},"current_reference":{"authors":["Dirk Bergemann","Juuso Välimäki"],"title":"Dynamic Mechanism Design: An Introduction","year":2018,"url":"https://cowles.yale.edu/sites/default/files/2022-09/d2102-r.pdf","doi":"10.1257/jel.20180892","locator":"§8 Concluding Remarks, pp.49–51 of June19,2018 revision","evidence_summary":"The conclusion explicitly identifies weaker common-prior assumptions, sequential participation, outside resale, and limited adoption of long-term contracts as research questions."},"formalization_note":"This is a finite-wealth, costless-exit adaptation of the sequential-participation agenda, not the already addressed unrestricted-bonding case."},"statusQualification":"Published question or research agenda verified at the cited locator. The mathematical specification is adapted for this catalogue and includes additional assumptions; source publication does not establish first-ever priority or certify this exact model remains unsolved. This is a finite-wealth, costless-exit adaptation of the sequential-participation agenda, not the already addressed unrestricted-bonding case.","statusReviewedAt":"2026-10-08"}
% FIRST_POSTED: 2026-10-08
% ENTRY_KIND: statement_only
% !TeX program = lualatex
\documentclass[11pt]{article}
\usepackage{fontspec}
\usepackage[a4paper,margin=25mm]{geometry}
\usepackage{amsmath,amssymb,mathtools}
\usepackage{xurl}
\usepackage[hidelinks]{hyperref}
\newcommand{\sref}[2]{\hyperlink{source:#1:#2}{[S#2]}}
\setlength{\emergencystretch}{3em}
\begin{document}
% problemId:30007573
\section*{Voluntary participation over time}\label{30007573}
\subsection{Definitions and mathematical statement}
Let one seller repeatedly allocate a service to \(n\) agents over \(T\) periods. Agent \(i\)'s type follows a known finite Markov chain and yields quasilinear discounted utility, with common discount factor \(\beta\in(0,1]\). A direct mechanism chooses feasible allocation \(x_t\) and nonnegative payment \(p_{it}\) from report histories. Agent \(i\) has observable initial liquid wealth \(B_i\), cannot borrow, and may exit after any private history with no future allocations or payments. Impose the pathwise constraint \(\sum_{s\leq t}p_{is}\leq B_i\), dynamic incentive compatibility against every reporting-and-exit strategy, and sequential participation
\[
 E\left[\sum_{s=t}^{T}\beta^{s-t}
 (v_i(x_s,\theta_{is})-p_{is})\mid h_{it}\right]\geq0
\]
at every truthful private history \(h_{it}\). Characterize revenue-optimal mechanisms and the loss relative to mechanisms requiring only initial participation, as functions of wealth, type persistence, and horizon. Determine whether limited refundable escrow or capped deposits recover part of that loss, stating the cash and exit rules exactly before comparing mechanisms. An existence or characterization result must include off-path incentive restrictions and histories reached after other agents' reports. The survey explicitly identifies sequential participation as less understood than dynamic truthfulness and notes that unrestricted initial bonding changes the problem; the finite-wealth specification here isolates the economically relevant limitation on that device.

\par\textbf{Formulation and scope.} This is a finite-wealth, costless-exit adaptation of the sequential-participation agenda, not the already addressed unrestricted-bonding case.

\par\textbf{Open-question provenance.} An explicit question or research agenda was verified at the source locator. This does not by itself certify that every later version remains unresolved \sref{30007573}{1}.

\par\textbf{Historical priority.} This is the earliest source in which the relevant agenda was verified during this audit; absolute priority has not been established.

\subsection{Short English statement}
Characterize optimal dynamic contracts under costless sequential exit and finite liquidity, and quantify how participation constraints or limited escrow reduce revenue.

\subsection{Sources}
\begin{itemize}\raggedright
\item[\textnormal{[S1]}]\hypertarget{source:30007573:1}{} Dirk Bergemann, Juuso Välimäki. 2018. Dynamic Mechanism Design: An Introduction. §8 Concluding Remarks, pp.49–51 of June19,2018 revision.
\url{https://cowles.yale.edu/sites/default/files/2022-09/d2102-r.pdf}.
DOI: 10.1257/jel.20180892.
{\small\textit{Earliest verified question/agenda reference; historical priority qualified below. The conclusion explicitly identifies weaker common-prior assumptions, sequential participation, outside resale, and limited adoption of long-term contracts as research questions.}}
\end{itemize}
\end{document}
